\chapter{Deploying Agents with AWS AgentCore}

Chapter~8 assembled the same assistant in two frameworks, producing two agent
implementations that run locally when their environment variables point to real gateway
and memory resources. Deployment extends that work into a governed cloud service: it
supplies an invocation boundary, an execution role, a memory resource, a governed tool
surface, a release mechanism, and the observability needed to diagnose failures in
production.

Amazon Bedrock AgentCore provides managed infrastructure for each of these concerns. AWS
CDK declares all required resources in code. Examining the CDK stack component by
component reveals what each AgentCore service does and why the provisioning order follows
the dependencies it does.

\section{Step 1: Understand what the runtime expects}

Before writing infrastructure, it helps to understand the contract the agent code already
satisfies. In chapter~8, both agents wrap themselves in a \texttt{BedrockAgentCoreApp}
and expose exactly one \texttt{@app.entrypoint} function:

\begin{lstlisting}[language=Python]
from bedrock_agentcore import BedrockAgentCoreApp

app = BedrockAgentCoreApp()

@app.entrypoint
async def invoke(payload: dict, context=None):
    prompt = str(payload.get("prompt") or "").strip()
    runtime_session_id = _runtime_session_id_from(
        payload, context
    )
    # ... agent logic ...
    yield {
        "message": {
            "role": "assistant",
            "content": [{"text": response_text}],
        }
    }

def main():
    app.run(host="0.0.0.0", port=8080)
\end{lstlisting}

\texttt{BedrockAgentCoreApp.run()} starts an HTTP server on port 8080. AgentCore Runtime
sends invocation requests to this server, supplies the \texttt{runtimeSessionId} through
the context object, and streams or collects the yielded events. The agent code does not
need to know it is running in a container or behind a managed runtime; it only needs to
be packaged so that AgentCore can start it.

This separation means the same code can be tested locally by running
\texttt{python~-m~agent.app} and passing test payloads directly, and deployed by
packaging the \texttt{agent/} directory as a Docker image.

\section{Step 2: Build the tool layer}

The runtime does not own the business logic. A well-structured deployment separates tool
implementation from agent orchestration by placing domain operations in a Lambda
function. The Lambda function is the first resource created in the CDK stack, before the
runtime, because the gateway and the runtime both depend on it.

\begin{lstlisting}
import { createToolsResources } from './tools-construct';

const tools = createToolsResources(this, {
  stageName: props.stage.name,
  removalPolicy,
  modelId: DEFAULT_MODEL_ID,
});
\end{lstlisting}

\texttt{createToolsResources} returns a Lambda function and a DynamoDB state table. The
Lambda function handles tool invocations from the AgentCore Gateway and any scheduled or
event-driven triggers the application requires. The DynamoDB table stores all typed
records with a configurable TTL.

This separation gives the deployment a clean ownership model: the agent owns reasoning,
the gateway owns tool discovery and protocol, the Lambda owns domain operations, and
DynamoDB owns persisted state.

\section{Step 3: Create the execution role}

The agent runtime needs an IAM role that grants it permission to invoke the model, write
logs, call the gateway, and access memory. AgentCore assumes this role using a service
principal with conditions that prevent confused deputy attacks:

\begin{lstlisting}
const runtimeRole = new iam.Role(this, 'AgentRuntimeRole', {
  assumedBy: new iam.ServicePrincipal(
    'bedrock-agentcore.amazonaws.com'
  ).withConditions({
    StringEquals: {
      'aws:SourceAccount': this.account,
    },
    ArnLike: {
      'aws:SourceArn':
        `arn:aws:bedrock-agentcore:${this.region}:`
        + `${this.account}:runtime/*`,
    },
  }),
});

runtimeRole.addToPolicy(new iam.PolicyStatement({
  actions: [
    'bedrock:InvokeModel',
    'bedrock:InvokeModelWithResponseStream',
    'cloudwatch:PutMetricData',
    'logs:CreateLogGroup',
    'logs:CreateLogStream',
    'logs:PutLogEvents',
  ],
  resources: ['*'],
}));

tools.toolFunction.grantInvoke(runtimeRole);
dataBucket.grantReadWrite(runtimeRole);
\end{lstlisting}

The conditions on the trust policy limit which AgentCore resources can assume this role.
Without them, any AgentCore runtime in the account could assume the role. The
\texttt{aws:SourceArn} condition restricts assumption to runtimes in this account and
region.

\section{Step 4: Provision AgentCore Memory}

Memory is provisioned as an independent resource before the runtime. Both the Strands and
LangGraph implementations receive the memory ID through an environment variable and bind
it at startup; they do not create memory resources themselves.

\begin{lstlisting}
import { Memory, MemoryStrategy } from
  '@aws-cdk/aws-bedrock-agentcore-alpha';

const memory = new Memory(this, 'AgentCoreMemory', {
  memoryName: `agent_memory_${props.stage.name}`,
  description: 'Long-term memory for user preferences.',
  expirationDuration: Duration.days(30),
  memoryStrategies: [
    MemoryStrategy.usingUserPreference({
      name: 'user_preferences',
      description: 'Long-term preferences across sessions.',
      namespaces: ['/actors/{actorId}/preferences'],
    }),
  ],
});

memory.grantWrite(runtimeRole);
memory.grantRead(runtimeRole);
\end{lstlisting}

The memory strategy declares that this resource stores user preferences keyed by actor
identifier. AgentCore Memory handles the vector index and retrieval. The agent code uses
the memory ID to construct the checkpointer, the store, or the session manager depending
on the framework, and the namespace \texttt{/actors/\{actorId\}/preferences} is where
both the Strands session manager and the LangGraph pre-model hook write and search.

\section{Step 5: Create the AgentCore Gateway}

The gateway is the MCP endpoint that the agent uses at runtime. It advertises tool
schemas, enforces IAM authentication on every tool invocation, and routes calls to the
Lambda backend.

\begin{lstlisting}
import {
  Gateway,
  GatewayAuthorizer,
  GatewayProtocol,
  MCPProtocolVersion,
  McpGatewaySearchType,
} from '@aws-cdk/aws-bedrock-agentcore-alpha';

const agentCoreGateway = new Gateway(scope,
  'AgentGateway', {
  gatewayName: `agent-gateway-${stageName}`,
  description: 'MCP gateway for agent tools.',
  authorizerConfiguration: GatewayAuthorizer.usingAwsIam(),
  protocolConfiguration: GatewayProtocol.mcp({
    supportedVersions: [
      MCPProtocolVersion.MCP_2025_06_18,
      MCPProtocolVersion.MCP_2025_03_26,
    ],
    searchType: McpGatewaySearchType.SEMANTIC,
    instructions:
      'Use these tools to retrieve data, trigger processing '
      + 'jobs, inspect results, and record governed decisions.',
  }),
});

agentCoreGateway.addLambdaTarget(
  'AgentToolsTarget', {
  gatewayTargetName: 'agent-tools',
  description: 'Domain tools via Lambda.',
  lambdaFunction: toolFunction,
  toolSchema: ToolSchema.fromInline(AGENT_TOOL_SCHEMA),
});

agentCoreGateway.grantInvoke(runtimeRole);
\end{lstlisting}

The gateway is not the Lambda function. The Lambda function is a target: it receives
invocations from the gateway when a tool is called. The gateway owns the MCP protocol
layer, handles tool discovery, and enforces SigV4 authentication. This means the agent
code calls the gateway endpoint, not the Lambda directly. The gateway decouples the agent
from the underlying implementation so a tool target can be a Lambda, a Fargate service,
or any HTTP endpoint without changing the agent code.

The \texttt{grantInvoke(runtimeRole)} call allows the runtime container to call the
gateway. Without it, the signed requests from the agent would fail with an authorisation
error.

The runtime role also needs a specific gateway invocation permission:

\begin{lstlisting}
runtimeRole.addToPolicy(new iam.PolicyStatement({
  actions: ['bedrock-agentcore:InvokeGateway'],
  resources: [
    agentCoreGateway.agentCoreGatewayArn,
    `${agentCoreGateway.agentCoreGatewayArn}`
    + `/gateway-endpoint/*`,
  ],
}));
\end{lstlisting}

\section{Step 6: Package and deploy the Runtime}

With the role, memory, and gateway in place, the runtime can be declared. The CDK
\texttt{Runtime} construct packages the \texttt{agent/} directory as a Docker image and
registers it with AgentCore:

\begin{lstlisting}
import {
  AgentRuntimeArtifact,
  Runtime,
  RuntimeAuthorizerConfiguration,
  RuntimeNetworkConfiguration,
  ProtocolType,
} from '@aws-cdk/aws-bedrock-agentcore-alpha';
import { ecrAssets } from 'aws-cdk-lib';

const runtime = new Runtime(this, 'AgentRuntime', {
  runtimeName: props.stage.runtimeName,
  executionRole: runtimeRole,
  agentRuntimeArtifact: AgentRuntimeArtifact.fromAsset(
    path.join(__dirname, '..', 'agent'), {
      platform: ecrAssets.Platform.LINUX_ARM64,
    }
  ),
  networkConfiguration:
    RuntimeNetworkConfiguration.usingPublicNetwork(),
  protocolConfiguration: ProtocolType.HTTP,
  authorizerConfiguration:
    RuntimeAuthorizerConfiguration.usingIAM(),
  environmentVariables: {
    STAGE: props.stage.name,
    AWS_REGION: this.region,
    MODEL_ID: DEFAULT_MODEL_ID,
    GATEWAY_URL: agentCoreGateway.agentCoreGatewayUrl,
    GATEWAY_REGION: this.region,
    MEMORY_ID: memory.memoryId,
    MEMORY_REGION: this.region,
  },
});
\end{lstlisting}

Several details are worth examining.

\texttt{AgentRuntimeArtifact.fromAsset} builds a Docker image from the \texttt{agent/}
directory and pushes it to ECR. The agent directory must contain a \texttt{Dockerfile}
that installs dependencies and sets the container entry point to the agent application
module. CDK handles the build and push during \texttt{cdk deploy}.

\texttt{environmentVariables} is the bridge between infrastructure and code. The agent
reads the gateway URL, memory ID, and model identifier from environment variables. This
keeps the same container image portable across stages: a personal stage, a gamma stage,
and a production stage all use the same image but receive different environment variables
at runtime.

\texttt{RuntimeAuthorizerConfiguration.usingIAM()} means that every caller of the
runtime must present a valid SigV4 signature. Callers need the
\texttt{bedrock-agentcore:InvokeAgentRuntime} permission for the runtime ARN.

\section{Step 7: Add the proxy Lambda for browser clients}

Browser applications cannot sign requests with SigV4 directly. The backend solves this
with a runtime proxy Lambda that sits between the API Gateway and the AgentCore runtime:

\begin{lstlisting}
const runtimeProxyFn = new lambda.Function(this,
  'RuntimeProxyFunction', {
  runtime: lambda.Runtime.PYTHON_3_11,
  handler: 'index.handler',
  code: lambda.Code.fromAsset(
    path.join(__dirname, '..', 'lambda', 'runtime-proxy')
  ),
  timeout: Duration.seconds(30),
  memorySize: 512,
  environment: {
    ALLOWED_RUNTIME_PATTERN: `^${props.stage.runtimeName}`,
    CORS_ALLOWED_ORIGINS: '*',
  },
});

runtimeProxyFn.addToRolePolicy(new iam.PolicyStatement({
  actions: ['bedrock-agentcore:InvokeAgentRuntime'],
  resources: [
    runtime.agentRuntimeArn,
    `${runtime.agentRuntimeArn}/runtime-endpoint/*`,
  ],
}));

const invokeResource = api.root
  .addResource('runtime')
  .addResource('invoke');

invokeResource.addMethod('POST',
  new apigateway.LambdaIntegration(runtimeProxyFn), {
    authorizationType: apigateway.AuthorizationType.IAM,
  }
);
\end{lstlisting}

The proxy Lambda accepts a request from the IAM-protected API Gateway endpoint, validates
the runtime name against the \texttt{ALLOWED\_RUNTIME\_PATTERN} environment variable,
constructs the \texttt{runtimeSessionId} from the request, calls
\texttt{bedrock-agentcore:InvokeAgentRuntime}, parses the streaming response, and returns
a JSON envelope to the caller.

The \texttt{ALLOWED\_RUNTIME\_PATTERN} check prevents callers from invoking arbitrary
runtimes by constructing a different runtime name. The API Gateway endpoint is protected
by IAM, so the browser application authenticates with IAM credentials (typically via
Amazon Cognito) and the proxy does the SigV4 signing on its behalf.

This proxy pattern is not convenience wrapping. It is a policy enforcement point: it
controls which runtimes can be invoked, how session identifiers are constructed, and how
streaming output is reduced to a response that the browser can parse.

\section{Step 8: Deploy and verify}

The stack is deployed with the CDK CLI:

\begin{lstlisting}
# Install CDK dependencies
npm install

# Deploy to a personal stage
npm run deploy:personal

# Or use CDK directly
npx cdk deploy AgentBackendStack-personal
\end{lstlisting}

After deployment, CDK outputs the runtime ARN, the API URL, the gateway URL, and the
memory ID. These values are the concrete links between the CDK infrastructure and the
agent's environment variables.

A minimal smoke test verifies the runtime is alive without invoking any backend tools or
the model:

\begin{lstlisting}[language=Python]
import boto3, json, uuid

def invoke_agent(
    runtime_arn: str,
    region: str,
    prompt: str,
) -> str:
    client = boto3.client(
        "bedrock-agentcore", region_name=region
    )
    response = client.invoke_agent_runtime(
        payload=json.dumps(
            {"prompt": prompt}
        ).encode("utf-8"),
        runtimeSessionId=f"testuser___{uuid.uuid4()}",
        agentRuntimeArn=runtime_arn,
    )
    body = response.get("response") or response.get("body")
    return extract_response_text(iter_objects(body))

# Send a health prompt and verify the container is alive
text = invoke_agent(runtime_arn, region, prompt="health")
print(text)  # expected: {"status": "ok", ...}
\end{lstlisting}

The health prompt takes the fast path in both agent implementations: it skips the model,
the gateway tools, and the memory layer, and returns a static status object. This makes
it safe to call frequently as a liveness check without incurring model costs. For a
fuller integration test, send a real task prompt and use Langfuse to capture the trace.
The evaluation chapter treats that trace as the foundation of diagnostic scoring.

\section{What the runtime owns versus what it delegates}

The deployment creates a clean separation of ownership across services.

\begin{longtable}{p{0.28\textwidth}p{0.62\textwidth}}
\toprule
\textbf{Service} & \textbf{Responsibility} \\
\midrule
AgentCore Runtime
    & Packages and runs the agent container; manages invocation
      lifecycle and session routing \\
\addlinespace
AgentCore Gateway
    & Exposes backend tools via MCP; enforces SigV4
      authentication on every tool call; handles protocol
      versioning \\
\addlinespace
AgentCore Memory
    & Stores conversational state, checkpoints, and long-term
      preferences; provides semantic retrieval \\
\addlinespace
Lambda (tools)
    & Implements all domain operations; owns business logic
      and state persistence \\
\addlinespace
DynamoDB
    & Persists typed records with a configurable TTL \\
\addlinespace
EventBridge
    & Triggers scheduled or event-driven processing pipelines \\
\addlinespace
API Gateway
    & Provides browser-facing ingress with IAM authentication
      and CORS \\
\addlinespace
Lambda (proxy)
    & Enforces runtime name policy; translates between API
      Gateway and AgentCore invocation protocol \\
\addlinespace
CloudWatch
    & Stores structured logs from all Lambda functions; enables
      query-based diagnostics \\
\addlinespace
Secrets Manager
    & Stores runtime secrets so they are not baked into the
      container image \\
\bottomrule
\end{longtable}

The model is surrounded by conventional cloud services. It does not touch infrastructure
directly. When the agent calls a gateway tool, IAM governs the call, the gateway
validates the schema, the Lambda runs the domain logic, and DynamoDB stores the result.
The agent sees a tool response.

\section{Strands and LangGraph in the same runtime}

Both frameworks deploy to AgentCore using the same CDK stack. The only structural
difference is which Python module in the \texttt{agent/} directory is set as the entry
point in the Dockerfile. The CDK runtime artifact, the IAM role, the gateway, and the
memory resource are identical in both cases.

This is the key property of the deployment pattern: it is framework-neutral. The runtime
does not care whether the agent is a Strands \texttt{Agent} object or a LangGraph
\texttt{create\_react\_agent} graph. It cares that the agent exposes an HTTP server on
port 8080, accepts a JSON payload, yields events in the expected format, and responds to
health checks.

The same memory ID works with the Strands \texttt{AgentCoreMemorySessionManager} and
with the LangGraph \texttt{AgentCoreMemorySaver} and \texttt{AgentCoreMemoryStore}. The
same gateway URL works with the Strands \texttt{GatewayToolPool} and with the LangGraph
\texttt{MultiServerMCPClient}. The framework choice is an application-layer decision; the
deployment layer does not need to change when that decision changes.

\section{Summary}

Deploying an agent to AgentCore follows a fixed order: tool layer first, then IAM role,
then memory, then gateway, then runtime, then proxy. Each resource depends on resources
created before it. The CDK stack encodes this dependency graph as TypeScript and
\texttt{cdk deploy} resolves it into a set of CloudFormation change sets.

The agent code stays portable throughout. It reads its gateway URL, memory ID, and model
identifier from environment variables. The same code runs locally against a test gateway,
on a personal stage, and in production. The deployment layer supplies environment
variables and handles authentication; the agent code handles reasoning. That separation
is what allows the same assistant to run as a Strands agent or a LangGraph graph without
modifying the infrastructure at all.